\documentclass[11pt,a4paper]{article}
\usepackage[utf8]{inputenc}
\usepackage[margin=1in]{geometry}
\usepackage{graphicx}
\usepackage{booktabs}
\usepackage{hyperref}
\usepackage{xcolor}
\usepackage{float}
\usepackage{fancyhdr}
\usepackage{titlesec}

% Header/Footer
\pagestyle{fancy}
\fancyhf{}
\rhead{Assignment 3 - MLflow Tracking}
\lhead{MLOps Course}
\rfoot{Page \thepage}

% Colors
\definecolor{winnergreen}{RGB}{34, 139, 34}
\definecolor{warningred}{RGB}{178, 34, 34}

\title{\textbf{Assignment 3: Observable ML with MLflow}\\[0.5em]
\large GAN Training Experiment Tracking \& Analysis}
\author{Student A}
\date{March 12, 2026}

\begin{document}

\maketitle

\section{Introduction}
This report documents the MLflow experiment tracking for a Vanilla GAN trained on the Fashion-MNIST dataset. Five experiments were conducted with varying hyperparameters to analyze the effect of learning rate and batch size on model convergence and stability.

\section{Experiment Configuration}

\begin{table}[H]
\centering
\caption{Hyperparameter Configurations for 5 Runs}
\begin{tabular}{@{}ccccc@{}}
\toprule
\textbf{Run} & \textbf{Learning Rate} & \textbf{Batch Size} & \textbf{Epochs} & \textbf{Latent Dim} \\
\midrule
1 & 0.0002 & 512 & 5 & 100 \\
2 & 0.001  & 512 & 5 & 100 \\
3 & 0.0001 & 512 & 5 & 100 \\
4 & 0.0002 & 256 & 5 & 100 \\
5 & 0.0002 & 128 & 5 & 100 \\
\bottomrule
\end{tabular}
\end{table}

\section{MLflow Screenshots}

\subsection{Table View - All 5 Runs Sorted by Accuracy}

\begin{figure}[H]
\centering
\includegraphics[width=0.9\textwidth]{image.png}
\caption{MLflow Table View showing all 5 experiment runs sorted by Discriminator Accuracy}
\label{fig:table_view}
\end{figure}

\subsection{Comparison View - Loss/Accuracy Curves}

\begin{figure}[H]
\centering
\includegraphics[width=0.9\textwidth]{runs compare.png}
\caption{MLflow Comparison View showing training curves for multiple runs}
\label{fig:comparison_view}
\end{figure}

\subsection{Artifact Gallery - Saved Model \& Environment Files}

\begin{figure}[H]
\centering
\includegraphics[width=0.9\textwidth]{artifact.png}
\caption{MLflow Artifacts tab showing saved Generator model and environment files}
\label{fig:artifacts}
\end{figure}

\section{Results Summary}

\begin{table}[H]
\centering
\caption{Final Metrics for All Runs}
\begin{tabular}{@{}ccccl@{}}
\toprule
\textbf{Run} & \textbf{Final D\_Loss} & \textbf{Final G\_Loss} & \textbf{Final D\_Acc} & \textbf{Notes} \\
\midrule
\textcolor{winnergreen}{\textbf{1}} & \textcolor{winnergreen}{\textbf{0.1238}} & \textcolor{winnergreen}{\textbf{3.70}} & \textcolor{winnergreen}{\textbf{99.5\%}} & \textcolor{winnergreen}{\textbf{Winner}} \\
2 & 0.0051 & 8.89 & 99.5\% & Unstable G\_loss \\
3 & 0.0644 & 3.94 & 98.9\% & Slower convergence \\
4 & 0.0526 & 6.69 & 99.6\% & Good accuracy \\
5 & 0.0989 & 6.94 & 98.9\% & More variance \\
\bottomrule
\end{tabular}
\end{table}

\section{Analysis}

\subsection{Winner: Run 1 (lr=0.0002, batch\_size=512)}

\textbf{Run 1} is the clear winner for the following reasons:

\begin{enumerate}
    \item \textbf{Lowest Generator Loss (3.70):} Among all runs, Run 1 achieved the lowest final generator loss, indicating the generator learned to produce images that effectively fool the discriminator.
    
    \item \textbf{Balanced GAN Equilibrium:} The discriminator loss (0.1238) and accuracy (99.5\%) suggest a healthy adversarial balance. The discriminator is neither too strong (which would starve the generator) nor too weak (which would lead to mode collapse).
    
    \item \textbf{Stable Training:} Unlike Run 2 (high learning rate), Run 1 showed consistent, monotonic improvement without erratic spikes in the loss curves.
\end{enumerate}

\subsection{Evidence of Training Issues}

\subsubsection{Run 2 - Unstable Training (High Learning Rate)}
\textcolor{warningred}{\textbf{Warning:}} Run 2 (lr=0.001) exhibited severe instability:
\begin{itemize}
    \item Generator loss spiked to \textbf{67.8} during Epoch 2 before recovering
    \item Final G\_loss of 8.89 is more than 2x higher than Run 1
    \item This is classic evidence of \textbf{learning rate being too high} for GAN training
\end{itemize}

\subsubsection{Run 3 - Slower Convergence (Low Learning Rate)}
Run 3 (lr=0.0001) showed:
\begin{itemize}
    \item More gradual loss decrease (smoother curves)
    \item Slightly lower final accuracy (98.9\%)
    \item Would likely improve with more epochs, but converges slower
\end{itemize}

\subsection{Overfitting/Underfitting Analysis}

\begin{itemize}
    \item \textbf{No clear overfitting observed:} GANs don't have a traditional train/validation split, but the discriminator accuracies plateauing near 99\% without diverging losses suggest healthy training.
    
    \item \textbf{Potential mode collapse indicators:} Very high discriminator accuracy (99.5\%+) combined with high generator loss could indicate mode collapse. Run 2's pattern (D\_acc=99.5\%, G\_loss=8.89) warrants investigation of generated sample diversity.
    
    \item \textbf{Underfitting in Run 3:} The lower learning rate resulted in the model not fully converging within 5 epochs, suggesting underfitting that could be resolved with longer training.
\end{itemize}

\section{Conclusion}

MLflow proved invaluable for tracking and comparing GAN experiments. The baseline configuration (lr=0.0002, batch\_size=512) emerged as optimal, demonstrating that:

\begin{itemize}
    \item GAN training is highly sensitive to learning rate (Run 2's instability)
    \item Batch size had less impact than learning rate in this experiment
    \item Observable ML practices enable data-driven hyperparameter decisions
\end{itemize}

\vspace{1em}
\noindent\textbf{Experiment tracked at:} \url{http://localhost:5000/#/experiments/1}

\end{document}
